\section{Локальные данные и сканер}

\subsection{Хранение и загрузка данных}

\begin{tabularx}{\textwidth}{L{48mm}Y}
\toprule
\textbf{Вопрос} & \textbf{Решение} \\
\midrule
Где хранятся данные? & Каталог товаров, ячейки, задания, результаты операций и отчёты о приёмке хранятся только в локальной базе данных устройства. \\
Как данные попадают в приложение? & Перед началом работы в приложение импортируются подготовленные файлы со справочником товаров, ячейками и заданиями. Внешняя учётная система не подключается. \\
Что содержит штрихкод? & Используется EAN-13. По отсканированному коду приложение находит конкретный товар в локальном каталоге и проверяет его соответствие открытому заданию. \\
Когда данные меняются? & Только после явного подтверждения сотрудником. Сканирование и ручной ввод только проверяют код. \\
Где хранится фото накладной? & Фотография подписанной накладной сохраняется в отчёте о приёмке и доступна при его просмотре. \\
\bottomrule
\end{tabularx}

\subsection{Поведение сканера}

\begin{tabularx}{\textwidth}{L{48mm}Y}
\toprule
\textbf{Ситуация} & \textbf{Решение} \\
\midrule
Основной способ сканирования & Камера Android-устройства; распознавание EAN-13 выполняется без сети. \\
Альтернативный способ & Сотрудник может вручную ввести 13 цифр EAN-13. Ручной ввод проходит такую же проверку, как код с камеры. \\
Успешное считывание & Камера приостанавливается, код проверяется в локальном каталоге и в контексте открытого задания. Затем сотрудник возвращается к подтверждению операции. \\
Ошибка считывания & Приложение показывает понятную причину ошибки, не изменяет локальные данные и позволяет повторить сканирование или перейти к ручному вводу. \\
Замена сканера & Бизнес-операции используют единый интерфейс сканера; камерная и аппаратная реализации не меняют сценарии экранов. \\
\bottomrule
\end{tabularx}
